\begin{center}\begin{tikzpicture}[font=\small,>=stealth,line join=round]
\draw[gray!40,rounded corners=3pt] (0,0) rectangle (14.2,8.3);\node at (7.1,7.700000000000001) {左右点不必同高，但连线的中点始终是中心};
\begin{scope}[shift={(1.491304347826087,2.4166666666666665)},x=1.3043478260869565cm,y=0.5111111111111111cm]\draw[->,gray] (-0.3,0)--(4.3,0) node[right]{$x$};\draw[->,gray] (0,-1.5)--(0,7.5) node[above]{$y$};\draw[blue,thick,solid,domain=0.4:3.6,samples=180] plot(\x,{(\x-2)^3+3});\fill[orange!85!black] (1,2) circle(2.2pt);\fill[orange!85!black] (3,4) circle(2.2pt);\draw[orange!85!black] (1,2)--(3,4);\draw[gray,dashed] (2,-1.2)--(2,7.2);\draw[gray,dashed] (-.2,3)--(4.2,3);\fill[violet] (2,3) circle(2.6pt);\end{scope}\node[anchor=west,align=left,text width=5.0cm] at (8.4,4.0) {$f=(x-2)^3+3$\\[9pt]$P(3,4),\ Q(1,2)$\\[9pt]$\frac{3+1}{2}=2,\quad\frac{4+2}{2}=3$\\[9pt]中点 $C(2,3)$\\[9pt]$f(2+t)+f(2-t)=6$};
\node[text width=12.6cm,align=center] at (7.1,.65) {这条具体曲线示范一般公式；严格判定仍要求所有点都有对应点。};
\end{tikzpicture}\end{center}
